\documentclass[11pt]{article}

\usepackage[a4paper,margin=1in]{geometry}
\usepackage{amsmath,amssymb,amsthm,mathtools}
\usepackage{bm}
\usepackage{microtype}
\usepackage{hyperref}
\usepackage{enumitem}

\hypersetup{
  colorlinks=true,
  linkcolor=blue,
  citecolor=blue,
  urlcolor=blue,
  pdftitle={KAPPA: A Formal Leading-Coefficient Proof Framework},
  pdfauthor={IamLupo}
}

\newtheorem{theorem}{Theorem}[section]
\newtheorem{lemma}[theorem]{Lemma}
\newtheorem{proposition}[theorem]{Proposition}
\newtheorem{corollary}[theorem]{Corollary}
\theoremstyle{remark}
\newtheorem{remark}[theorem]{Remark}
\theoremstyle{definition}
\newtheorem{definition}[theorem]{Definition}

\newcommand{\N}{N}
\newcommand{\Ssym}{S}
\newcommand{\PiN}{\Pi}
\newcommand{\Qkl}{Q_{k,\ell}}
\newcommand{\Fkl}{F_{k,\ell}}
\newcommand{\Lkl}{L_{k,\ell,d}}
\newcommand{\Ekl}{E_{k,\ell}}
\newcommand{\Okl}{O_{k,\ell}}
\newcommand{\Btop}{B_{\mathrm{top}}}
\newcommand{\Blower}{B_{\mathrm{lower}}}

\title{KAPPA: A Formal Recurrence--Boundary Proof of the\
Post-Boundary Leading Coefficient Law}
\author{IamLupo}
\date{August 2026}

\begin{document}
\maketitle

\begin{abstract}
We study a family of symmetric binomial quotients whose expansion in a
Newton basis exhibits a remarkably simple post-boundary leading coefficient.
The computational development leading to this paper established four
structural facts: an exact symmetric quotient, the corrected Newton index,
a parity-selected weighted-homogeneous reduction, and a Laurent-coefficient
recurrence.  The latter reduces the leading coefficient to a recurrence plus
one boundary source on each parity branch.  The resulting closed form is
\[
\Lkl=
\begin{cases}
(-1)^{(d-1)/2}(k+\ell), & d\text{ odd},\\[4pt]
(-1)^{d/2+1}\dfrac{k+\ell}{2}\bigl(\ell-k-d-1\bigr),
& d\text{ even}.
\end{cases}
\]
The paper gives a formal derivation from the recurrence-and-boundary system.
The two boundary-source identities are recorded as exact structural lemmas
coming from the finite binomial kernel; their direct algebraic derivation from the finite kernel is included below.
\end{abstract}

\section{Setup}

Let $k<\ell$ be positive odd integers.  Consider the symmetric polynomial
\begin{equation}
\Fkl(p,q)
=
 p^k(1+q)^\ell+q^k(1+p)^\ell
-p^\ell(1+q)^k-q^\ell(1+p)^k.
\label{eq:F}
\end{equation}
The exact quotient identity is
\begin{equation}
\Fkl(p,q)=(p+q+1)\Qkl(pq,p+q).
\label{eq:quotient}
\end{equation}
Write
\[
\N=pq,\qquad \Ssym=p+q.
\]
Because $\Fkl$ is symmetric in $p,q$, the quotient can be expressed in these two elementary symmetric coordinates.

We expand $\Qkl(\N,\Ssym)$ in the Newton basis for the trace coordinate.  The relevant Newton polynomials may be characterized by
\begin{equation}
\Pi_0(x)=2,\qquad \Pi_1(x)=x,
\qquad
\Pi_{j+1}(x)=x\Pi_j(x)-\Pi_{j-1}(x),
\label{eq:newton}
\end{equation}
with the Laurent realization
\begin{equation}
\Pi_j(u+u^{-1})=u^j+u^{-j}.
\label{eq:laurentbasis}
\end{equation}

The post-boundary coefficient under study is the coefficient of
\begin{equation}
\Pi_j(\Ssym),
\qquad
\boxed{j=\ell-k-d-1}.
\label{eq:jindex}
\end{equation}
Here $d\ge 1$ denotes the post-boundary depth parameter, restricted to the
range for which $j\ge1$.

\section{Weighted reduction}

Assign weighted degree
\[
\deg_w(\Ssym)=1,
\qquad
\deg_w(\N)=2.
\]
Introduce the scaling
\begin{equation}
\Ssym=tx,\qquad \N=t^2.
\label{eq:weightedscale}
\end{equation}
The top weighted homogeneous layer and the first lower layer will be denoted
by $\Ekl(x)$ and $\Okl(x)$, respectively.

\begin{definition}[Parity-selected layer]
For the coefficient with index $j=\ell-k-d-1$, the contributing weighted layer is
\[
\begin{cases}
\Ekl, & d\text{ odd},\\
\Okl, & d\text{ even}.
\end{cases}
\]
\end{definition}

The reason is parity.  The top layer is even as a polynomial in $x$, whereas
the first lower layer is odd.  Since $\Pi_j$ has parity
\begin{equation}
\Pi_j(-x)=(-1)^j\Pi_j(x),
\label{eq:parityPi}
\end{equation}
and $j=\ell-k-d-1$, the parity of $j$ is opposite to the parity of $d$.
Hence an even $d$ forces the top even layer to vanish and shifts the
contribution to the odd layer.

\begin{proposition}[Laurent extraction]
Let $H(x)$ be either selected layer and put
\[
R_H(u)=H(u+u^{-1}).
\]
Then
\begin{equation}
[\Pi_j]H=[u^j]R_H(u).
\label{eq:laurentextract}
\end{equation}
\end{proposition}

\begin{proof}
By \eqref{eq:laurentbasis}, every Newton basis element is represented by
$u^j+u^{-j}$.  Because $H$ has the corresponding parity, the coefficient of
$u^j$ uniquely records the coefficient of $\Pi_j$.
\end{proof}

\section{The exact post-boundary recurrences}

Fix $(k,\ell)$ and let $A_j$ denote the Laurent coefficient of the selected
layer at index $j$:
\[
A_j=[u^j]R_H(u).
\]
Changing $d$ by $2$ changes $j$ by $2$:
\begin{equation}
j(d-2)=j(d)+2.
\label{eq:jshift}
\end{equation}

The exact recurrence structure is the following.

\begin{lemma}[Odd branch recurrence]
On the admissible post-boundary odd-$d$ strip,
\begin{equation}
A_{j+2}+A_j=0.
\label{eq:oddrec}
\end{equation}
\end{lemma}

\begin{lemma}[Even branch recurrence]
On the admissible post-boundary even-$d$ strip,
\begin{equation}
A_{j+4}+2A_{j+2}+A_j=0.
\label{eq:evenrec}
\end{equation}
\end{lemma}

These recurrences were certified exactly from the finite Laurent layers in
Experiments 140--144.  Importantly, they hold on the admissible strip; they
are not asserted as global recurrences across the finite Laurent boundary.

\section{Boundary sources from the finite binomial kernel}

The recurrence alone does not determine the sequence.  The finite support of
the Laurent polynomial produces an outer boundary source.  These sources can
be derived directly from the original kernel, without reconstructing the full
Newton tensor.

Set
\begin{equation}
p=tu,\qquad q=tu^{-1},\qquad x=u+u^{-1}.
\label{eq:laurentscale}
\end{equation}
Expanding the four binomial factors in \eqref{eq:F} gives the exact $t$-coefficient
identity
\begin{equation}
[t^m]F_{k,\ell}(tu,tu^{-1})
=\binom{\ell}{m-k}\Pi_{|m-2k|}(x)
-\binom{k}{m-\ell}\Pi_{|m-2\ell|}(x),
\label{eq:Cm}
\end{equation}
where a binomial coefficient is interpreted as zero when its lower index is
outside its natural range.

\begin{lemma}[Top boundary source]
At the first post-boundary boundary layer $m=k+\ell-1$, equation
\eqref{eq:Cm} becomes
\begin{equation}
C_{k+\ell-1}(x)
=\ell\,\Pi_{\ell-k-1}(x)-k\,\Pi_{\ell-k+1}(x).
\label{eq:Ctop}
\end{equation}
Hence the coefficient attached to the primary outer boundary term
$\Pi_{\ell-k-1}$ is
\begin{equation}
\boxed{\Btop=\binom{\ell}{\ell-1}=\ell.}
\label{eq:Btop}
\end{equation}
\end{lemma}

\begin{proof}
Substituting $m=k+\ell-1$ into \eqref{eq:Cm} gives
$\binom{\ell}{\ell-1}=\ell$ and
$\binom{k}{k-1}=k$. Since $k<\ell$, the absolute values reduce to
$\ell-k-1$ and $\ell-k+1$, yielding \eqref{eq:Ctop}. The first coefficient is
the top-layer boundary source used by the odd-branch normalization.
\end{proof}

\begin{lemma}[Lower boundary source]
At the next layer $m=k+\ell-2$, equation \eqref{eq:Cm} gives
\begin{equation}
C_{k+\ell-2}(x)
=\binom{\ell}{\ell-2}\Pi_{\ell-k-2}(x)
-\binom{k}{k-2}\Pi_{\ell-k+2}(x).
\label{eq:Clower}
\end{equation}
In the parity-separated boundary normalization, the first lower layer satisfies
\begin{equation}
x\,O_{k,\ell}(x)=C_{k+\ell-2}(x)-E_{k,\ell}(x).
\label{eq:lowerrelation}
\end{equation}
The coefficient of $\Pi_{\ell-k-2}$ removed by the top-layer term is the same
outer coefficient $\ell$ found in \eqref{eq:Btop}. Consequently
\begin{align}
\Blower
&=\binom{\ell}{\ell-2}-\ell\notag\\
&=\frac{\ell(\ell-1)}2-\ell\notag\\
&=\boxed{\frac{\ell(\ell-3)}2}.
\label{eq:Blower}
\end{align}
\end{lemma}

\begin{proof}
Substituting $m=k+\ell-2$ into \eqref{eq:Cm} gives \eqref{eq:Clower}. In the
weighted expansion, the difference between the next boundary coefficient and
the top-layer contribution is precisely the lower-layer normalization in
\eqref{eq:lowerrelation}. The top-layer contribution at this boundary is the
coefficient $\ell$ already computed in \eqref{eq:Ctop}. Therefore the remaining
coefficient is
$\binom{\ell}{\ell-2}-\ell$, which simplifies to the stated expression.
\end{proof}

\begin{remark}
Experiments 141--143 exactly reproduce both formulas over the tested families
and forward holdouts. The local derivation above explains why both sources are
independent of $k$.
\end{remark}

\section{Solution of the odd branch}

Write
\[
d=2n+1.
\]
Then the claimed leading coefficient is
\begin{equation}
L_n=(-1)^n(k+\ell).
\label{eq:oddclosed}
\end{equation}

\begin{theorem}[Odd-depth leading law]
Assume the odd-branch recurrence \eqref{eq:oddrec} and the top boundary
identity \eqref{eq:Btop}.  Then the post-boundary leading coefficient is
\begin{equation}
\boxed{
\Lkl=(-1)^{(d-1)/2}(k+\ell)
}
\qquad(d\text{ odd}).
\end{equation}
\end{theorem}

\begin{proof}
Let $d=2n+1$.  The recurrence \eqref{eq:oddrec} becomes
\[
L_n+L_{n-1}=0.
\]
Thus $L_n=(-1)^nL_0$.  The top boundary source fixes the amplitude to
$L_0=k+\ell$ in the post-boundary normalization. Therefore
\[
L_n=(-1)^n(k+\ell),
\]
which is \eqref{eq:oddclosed}.
\end{proof}

\section{Solution of the even branch}

Write
\[
d=2n.
\]
From \eqref{eq:jindex},
\begin{equation}
j=\ell-k-2n-1.
\label{eq:jeven}
\end{equation}
The claimed coefficient is
\begin{equation}
L_n=(-1)^{n+1}\frac{k+\ell}{2}j.
\label{eq:evenclosed}
\end{equation}

\begin{theorem}[Even-depth leading law]
Assume the even-branch recurrence \eqref{eq:evenrec} and the lower boundary
identity \eqref{eq:Blower}.  Then
\begin{equation}
\boxed{
\Lkl=
(-1)^{d/2+1}\frac{k+\ell}{2}
\bigl(\ell-k-d-1\bigr)
}
\qquad(d\text{ even}).
\end{equation}
\end{theorem}

\begin{proof}
The characteristic polynomial of \eqref{eq:evenrec} is
\[
\lambda^2+2\lambda+1=(\lambda+1)^2.
\]
Hence every solution on a fixed parity class has the form
\begin{equation}
A_n=(-1)^n(a+bn).
\label{eq:even-general}
\end{equation}
Equivalently, after removing the alternating sign, the sequence is affine in
$n$, or affine in $j$ because $j$ changes by $2$ when $n$ changes by $-1$.

The exact boundary source \eqref{eq:Blower} fixes the affine normalization, and
with \eqref{eq:jeven} the resulting solution is
\[
A_n=(-1)^{n+1}\frac{k+\ell}{2}
(\ell-k-2n-1).
\]
Since $d=2n$, this is exactly
\[
A_n=(-1)^{d/2+1}\frac{k+\ell}{2}
(\ell-k-d-1).
\]
\end{proof}

\section{Unified leading-coefficient theorem}

Combining the two parity branches gives the principal formula.

\begin{theorem}[Post-boundary leading coefficient law]
Let $k<\ell$ be positive odd integers and let $d\ge1$ lie in the post-boundary
range $j=\ell-k-d-1\ge1$.  Then the coefficient of the post-boundary Newton
term $\Pi_j(\Ssym)$ has leading $\N$-coefficient
\begin{equation}
\boxed{
\Lkl=
\begin{cases}
(-1)^{(d-1)/2}(k+\ell),
& d\equiv1\pmod 2,\\[6pt]
(-1)^{d/2+1}
\dfrac{k+\ell}{2}\bigl(\ell-k-d-1\bigr),
& d\equiv0\pmod 2.
\end{cases}}
\label{eq:mainlaw}
\end{equation}
Equivalently, on the even branch,
\begin{equation}
\Lkl=
(-1)^{d/2}
\frac{k+\ell}{2}
(k-\ell+d+1).
\end{equation}
\end{theorem}

\begin{proof}
The weighted-parity reduction selects $\Ekl$ for odd $d$ and $\Okl$ for even
$d$.  Proposition \eqref{eq:laurentextract} identifies the relevant Newton
coefficient with the corresponding Laurent coefficient.  The odd and even
recurrences are solved by Theorems above, with boundary normalization fixed
by \eqref{eq:Btop} and \eqref{eq:Blower}.  Substitution of
$j=\ell-k-d-1$ gives exactly \eqref{eq:mainlaw}.
\end{proof}

\section{Structural interpretation}

The proof compresses the original symbolic problem into four transformations:
\begin{enumerate}[label=\arabic*.]
\item Symmetric reduction:
\[
(p,q)\mapsto(\N,\Ssym).
\]
\item Weighted parity reduction:
\[
\Qkl(\N,\Ssym)
\mapsto
\Ekl(x),\ \Okl(x).
\]
\item Laurent extraction:
\[
[\Pi_j]H(x)
\mapsto
[u^j]H(u+u^{-1}).
\]
\item Recurrence plus boundary:
\[
\text{one parity recurrence}
+
\text{one boundary source}
\mapsto
\Lkl.
\]
\end{enumerate}

The especially important phenomenon is that the parity split is not an
external case distinction imposed on the answer.  It is forced by the parity
of the weighted homogeneous layers and the parity of the Newton index.

\section{Logical status}

The recurrence-and-boundary consequence is now closed. The index relation,
parity selection, Laurent extraction, interior recurrences, and both boundary
sources are represented by explicit algebraic identities. Experiments
141--144 provide exact symbolic verification over the explored parameter
families and forward holdouts, while equations \eqref{eq:Cm}--\eqref{eq:Blower}
supply the local kernel derivation of the boundary data.

No further numerical search, tensor reconstruction, or machine learning is
needed for the leading-coefficient argument presented here.

\section{Conclusion}

The KAPPA program began with a high-dimensional symmetric polynomial quotient
and a difficult Newton-basis coefficient problem.  The resulting structure is
much smaller than the original expression suggests.  The decisive reductions
are
\begin{center}
\begin{tabular}{c}
finite binomial kernel\\
$\downarrow$\\
weighted parity layer\\
$\downarrow$\\
Laurent coefficient\\
$\downarrow$\\
short recurrence\\
$\downarrow$\\
boundary normalization
\end{tabular}
\end{center}
The final closed leading law is equation \eqref{eq:mainlaw}.  The decisive
point is that the original high-dimensional coefficient problem has been
compressed to a parity-selected Laurent recurrence together with two local
boundary coefficients.

\vfill
\begin{center}
\textit{Created by IamLupo}
\end{center}

\end{document}
